\chapter{MVC + Service Layer Architecture}

\section{Professional Backend Folder Structure}

\begin{diagrambox}[backend/src Folder Tree]
\begin{verbatim}
backend/src/
|-- config/db.js   controllers/   middleware/   models/
|-- routes/   services/   validators/   utils/
|-- app.js  (Express + middleware + routes)
`-- server.js  (HTTP listener + DB connection)
\end{verbatim}
\end{diagrambox}

\section{The Request Flow}

\begin{diagrambox}[Layered Request Flow]
\begin{verbatim}
Client -> Route (path match) -> Middleware (auth/validate)
       -> Controller (parse req, call service, shape res)
       -> Service (business rules) -> Model (Mongoose) -> MongoDB
\end{verbatim}
\end{diagrambox}

\section{Anti-Pattern: Everything in the Route Handler}

\begin{warnbox}[ANTI-PATTERN]
Mixing auth, validation, and DB access directly in a route handler does not scale past a handful of endpoints.
\end{warnbox}

\begin{lstlisting}[language=JavaScript]
// routes/task.routes.js  -- DO NOT DO THIS
router.post("/", async (req, res) => {
  const token = req.cookies.token;
  if (!token) return res.status(401).json({ message: "Not authenticated" });
  let userId;
  try { userId = jwt.verify(token, process.env.JWT_SECRET).id; }
  catch (e) { return res.status(401).json({ message: "Invalid token" }); }

  const { title, status } = req.body;
  if (!title?.trim()) return res.status(400).json({ message: "Title is required" });

  try {
    const count = await Task.countDocuments({ user: userId });
    if (count >= 100) return res.status(400).json({ message: "Task limit reached" });
    const task = new Task({ title, status, user: userId });
    await task.save();
    return res.status(201).json(await task.populate("user", "name email"));
  } catch (err) {
    return res.status(500).json({ message: err.message });
  }
});
\end{lstlisting}

\section{Refactored: Route -> Controller -> Service}

\subsection{routes/task.routes.js}

\begin{lstlisting}[language=JavaScript]
const { createTask } = require("../controllers/task.controller");
const { protect } = require("../middleware/auth.middleware");
const { validateTask } = require("../validators/task.validator");

router.post("/", protect, validateTask, createTask);
\end{lstlisting}

\subsection{controllers/task.controller.js}

\begin{lstlisting}[language=JavaScript]
const taskService = require("../services/task.service");

// Controller: translates HTTP <-> service calls. No Mongoose here.
const createTask = async (req, res, next) => {
  try {
    const task = await taskService.createTask(req.body, req.user.id);
    res.status(201).json(task);
  } catch (err) {
    next(err);
  }
};

module.exports = { createTask };
\end{lstlisting}

\subsection{services/task.service.js}

\begin{lstlisting}[language=JavaScript]
const Task = require("../models/Task");
const MAX_TASKS_PER_USER = 100;

// Service: owns business rules and all Mongoose interaction.
const createTask = async (taskData, userId) => {
  const count = await Task.countDocuments({ user: userId });
  if (count >= MAX_TASKS_PER_USER) {
    const err = new Error("Task limit reached");
    err.statusCode = 400;
    throw err;
  }
  const task = await Task.create({ ...taskData, user: userId });
  return task.populate("user", "name email");
};

module.exports = { createTask };
\end{lstlisting}

\section{Layer Responsibilities}

\begin{table}[h]
\centering
\begin{tabularx}{\textwidth}{l X X}
\toprule
\textbf{Layer} & \textbf{Responsibility} & \textbf{Should NOT contain} \\
\midrule
Route       & Map path to controller, attach middleware & Business logic, DB calls \\
Middleware  & Auth, validation, logging & Resource-specific rules \\
Controller  & Read req, call service, shape res & Direct queries, business rules \\
Service     & Business rules, model orchestration & req/res, status codes \\
Model       & Schema, field validation, hooks & Cross-resource rules \\
\bottomrule
\end{tabularx}
\end{table}

\begin{tipbox}[TIP]
If a function could run from a CLI script with no \texttt{req}/\texttt{res} in sight, it belongs in the service layer.
\end{tipbox}
